
\documentclass{article}
\usepackage{colortbl}
\usepackage{makecell}
\usepackage{multirow}
\usepackage{supertabular}

\begin{document}

\newcounter{utterance}

\centering \large Interaction Transcript for game `cladder', experiment `full\_v1.5\_default', episode 5573 with refineonly{-}Qwen3.5{-}9B.
\vspace{24pt}

{ \footnotesize  \setcounter{utterance}{1}
\setlength{\tabcolsep}{0pt}
\begin{supertabular}{c@{$\;$}|p{.15\linewidth}@{}p{.15\linewidth}p{.15\linewidth}p{.15\linewidth}p{.15\linewidth}p{.15\linewidth}}
    \# & \multicolumn{2}{c}{Player} && \multicolumn{2}{c}{Game Master} \\
    \hline

    \theutterance \stepcounter{utterance}  
    & & & \multicolumn{4}{p{0.6\linewidth}}{
        \cellcolor[rgb]{0.9,0.9,0.9}{
            \makecell[{{p{\linewidth}}}]{
                \texttt{\tiny{[P1$\langle$GM]}}
                \texttt{You are an expert in causal inference. The following question is not a typical commonsense query, but rather a meticulously designed question created by a professor specializing in causal inference, intended to assess the students' mastery of the course content.} \\
\texttt{Imagine a self{-}contained, hypothetical world with only the following conditions, and without any unmentioned factors or causal relationships: Gender has a direct effect on treatment and peanut allergy. Treatment has a direct effect on peanut allergy. For patients who are not male and not receiving treatment, the probability of being allergic to peanuts is 98\%. For patients who are not male and receiving treatment, the probability of being allergic to peanuts is 87\%. For patients who are male and not receiving treatment, the probability of being allergic to peanuts is 34\%. For patients who are male and receiving treatment, the probability of being allergic to peanuts is 14\%. The overall probability of male gender is 43\%. For patients receiving treatment, would it be less likely to see being allergic to peanuts if the patient had received no treatment?} \\
\texttt{Start your answer with 'Yes' or 'No', followed by additional reasoning or evidence to support your explanation.} \\
            }
        }
    }
     \\ \\

    \theutterance \stepcounter{utterance}  
    & \multicolumn{4}{p{0.6\linewidth}}{
        \cellcolor[rgb]{0.561,0.737,0.561}{
            \makecell[{{p{\linewidth}}}]{
                \texttt{\tiny{[P1$\rangle$GM]}}
                \texttt{No} \\
\\ 
\texttt{**Reasoning:**} \\
\\ 
\texttt{To determine if it would be less likely to see a patient allergic to peanuts if they had received no treatment, we must compare the probability of the outcome (peanut allergy) under the actual treatment condition (\$P(Y=1|T=1)\$) with the counterfactual probability under no treatment (\$P(Y=1|T=0)\$).} \\
\\ 
\texttt{Based on the data provided in the hypothetical world:} \\
\texttt{1.  **Actual Probability (\$P(Y=1|T=1)\$):** The problem states, "For patients receiving treatment, the probability of being allergic to peanuts is..." We must calculate the weighted average based on the gender distribution, as the treatment effect differs by gender.} \\
\texttt{    *   \$P(Male) = 0.43\$, so \$P(Not Male) = 1 {-} 0.43 = 0.57\$.} \\
\texttt{    *   For Not Male receiving treatment: \$P(Y=1|T=1, \textbackslash{}neg M) = 0.87\$.} \\
\texttt{    *   For Male receiving treatment: \$P(Y=1|T=1, M) = 0.14\$.} \\
\texttt{    *   \$P(Y=1|T=1) = (0.57 \textbackslash{}times 0.87) + (0.43 \textbackslash{}times 0.14) = 0.4959 + 0.0602 = 0.5561\$ (or 55.61\%).} \\
\\ 
\texttt{2.  **Counterfactual Probability (\$P(Y=1|T=0)\$):** This represents the probability of being allergic if the patient had *not* received treatment, holding their gender constant.} \\
\texttt{    *   For Not Male not receiving treatment: \$P(Y=1|T=0, \textbackslash{}neg M) = 0.98\$.} \\
\texttt{    *   For Male not receiving treatment: \$P(Y=1|T=0, M) = 0.34\$.} \\
\texttt{    *   \$P(Y=1|T=0) = (0.57 \textbackslash{}times 0.98) + (0.43 \textbackslash{}times 0.34) = 0.5586 + 0.1462 = 0.7048\$ (or 70.48\%).} \\
\\ 
\texttt{**Comparison:**} \\
\texttt{The probability of being allergic to peanuts if the patient had received no treatment (70.48\%) is **higher** than the probability if they received treatment (55.61\%).} \\
\\ 
\texttt{The question asks: "Would it be **less** likely to see being allergic to peanuts if the patient had received no treatment?"} \\
\texttt{Since \$P(Y=1|T=0) > P(Y=1|T=1)\$, it is actually *more* likely} \\
            }
        }
    }
    & & \\ \\

    \theutterance \stepcounter{utterance}  
    & & & \multicolumn{2}{p{0.3\linewidth}}{
        \cellcolor[rgb]{0.9,0.9,0.9}{
            \makecell[{{p{\linewidth}}}]{
                \texttt{\tiny{[GM$|$GM]}}
                \texttt{No} \\
            }
        }
    }
    & & \\ \\

    \theutterance \stepcounter{utterance}  
    & & & \multicolumn{2}{p{0.3\linewidth}}{
        \cellcolor[rgb]{0.9,0.9,0.9}{
            \makecell[{{p{\linewidth}}}]{
                \texttt{\tiny{[GM$|$GM]}}
                \texttt{no} \\
            }
        }
    }
    & & \\ \\

    \theutterance \stepcounter{utterance}  
    & & & \multicolumn{2}{p{0.3\linewidth}}{
        \cellcolor[rgb]{0.9,0.9,0.9}{
            \makecell[{{p{\linewidth}}}]{
                \texttt{\tiny{[GM$|$GM]}}
                \texttt{game\_result = WIN} \\
            }
        }
    }
    & & \\ \\

\end{supertabular}
}

\end{document}
